\subsubsection{\label{subsec.gf.defs.fps}The definition of formal power
series}

Until the end of Chapter \ref{chap.gf}, the following convention will be in place:

\begin{convention}
Fix a commutative ring $K$. (For example, $K$ can be $\mathbb{Z}$ or
$\mathbb{Q}$ or $\mathbb{C}$.)
\end{convention}

\begin{definition}
\label{def.fps.fps}A \emph{formal power series} (or, short, \emph{FPS}) in $1$
indeterminate over $K$ means a sequence $\left(  a_{0},a_{1},a_{2}%
,\ldots\right)  =\left(  a_{n}\right)  _{n\in\mathbb{N}}\in K^{\mathbb{N}}$ of
elements of $K$.
\end{definition}

Examples of FPSs over $\mathbb{Z}$ are $\left(  0,0,0,\ldots\right)  $ and
$\left(  1,0,0,0,\ldots\right)  $ and $\left(  1,1,1,1,\ldots\right)  $ and
$\left(  1,2,3,4,\ldots\right)  $.

Definition \ref{def.fps.fps} technically answers the question
\textquotedblleft what is an FPS\textquotedblright; however, the questions
\textquotedblleft what can we do with an FPS\textquotedblright\ and
\textquotedblleft why do the examples in Section \ref{sec.gf.exas}
work\textquotedblright\ or \textquotedblleft what is $x$\textquotedblright%
\ remain open. These questions will take us a while.

First, let us define some operations on FPSs.

\begin{definition}
\label{def.fps.ops}\textbf{(a)} The \emph{sum} of two FPSs $\mathbf{a}=\left(
a_{0},a_{1},a_{2},\ldots\right)  $ and $\mathbf{b}=\left(  b_{0},b_{1}%
,b_{2},\ldots\right)  $ is defined to be the FPS%
\[
\left(  a_{0}+b_{0},\ \ a_{1}+b_{1},\ \ a_{2}+b_{2},\ \ \ldots\right)  .
\]
It is denoted by $\mathbf{a}+\mathbf{b}$. \medskip

\textbf{(b)} The \emph{difference} of two FPSs $\mathbf{a}=\left(  a_{0}%
,a_{1},a_{2},\ldots\right)  $ and $\mathbf{b}=\left(  b_{0},b_{1},b_{2}%
,\ldots\right)  $ is defined to be the FPS%
\[
\left(  a_{0}-b_{0},\ \ a_{1}-b_{1},\ \ a_{2}-b_{2},\ \ \ldots\right)  .
\]
It is denoted by $\mathbf{a}-\mathbf{b}$. \medskip

\textbf{(c)} If $\lambda\in K$ and if $\mathbf{a}=\left(  a_{0},a_{1}%
,a_{2},\ldots\right)  $ is an FPS, then we define an FPS
\[
\lambda\mathbf{a}:=\left(  \lambda a_{0},\lambda a_{1},\lambda a_{2}%
,\ldots\right)  .
\]


\textbf{(d)} The \emph{product} of two FPSs $\mathbf{a}=\left(  a_{0}%
,a_{1},a_{2},\ldots\right)  $ and $\mathbf{b}=\left(  b_{0},b_{1},b_{2}%
,\ldots\right)  $ is defined to be the FPS $\left(  c_{0},c_{1},c_{2}%
,\ldots\right)  $, where
\begin{align*}
c_{n}  &  =\sum_{i=0}^{n}a_{i}b_{n-i}=\sum_{\substack{\left(  i,j\right)
\in\mathbb{N}^{2};\\i+j=n}}a_{i}b_{j}\\
&  =a_{0}b_{n}+a_{1}b_{n-1}+a_{2}b_{n-2}+\cdots+a_{n}b_{0}%
\ \ \ \ \ \ \ \ \ \ \text{for each }n\in\mathbb{N}.
\end{align*}
This product is denoted by $\mathbf{a}\cdot\mathbf{b}$ or just by
$\mathbf{ab}$. \medskip

\textbf{(e)} For each $a\in K$, we define $\underline{a}$ to be the FPS
$\left(  a,0,0,0,\ldots\right)  $. An FPS of the form $\underline{a}$ for some
$a\in K$ (that is, an FPS $\left(  a_{0},a_{1},a_{2},\ldots\right)  $
satisfying $a_{1}=a_{2}=a_{3}=\cdots=0$) is said to be \emph{constant}.
\medskip

\textbf{(f)} The set of all FPSs (in $1$ indeterminate over $K$) is denoted
$K\left[  \left[  x\right]  \right]  $.
\end{definition}

The following theorem is crucial: essentially it says that the operations on
FPSs that we have just defined behave as such operations should:

\begin{theorem}
\label{thm.fps.ring}\textbf{(a)} The set $K\left[  \left[  x\right]  \right]
$ is a commutative ring (with its operations $+$, $-$ and $\cdot$ defined in
Definition \ref{def.fps.ops}). Its zero and its unity are the FPSs
$\underline{0}=\left(  0,0,0,\ldots\right)  $ and $\underline{1}=\left(
1,0,0,0,\ldots\right)  $. This means, concretely, that the following facts hold:

\begin{enumerate}
\item \emph{Commutativity of addition:} We have $\mathbf{a}+\mathbf{b}%
=\mathbf{b}+\mathbf{a}$ for all $\mathbf{a},\mathbf{b}\in K\left[  \left[
x\right]  \right]  $.

\item \emph{Associativity of addition:} We have $\mathbf{a}+\left(
\mathbf{b}+\mathbf{c}\right)  =\left(  \mathbf{a}+\mathbf{b}\right)
+\mathbf{c}$ for all $\mathbf{a},\mathbf{b},\mathbf{c}\in K\left[  \left[
x\right]  \right]  $.

\item \emph{Neutrality of zero:} We have $\mathbf{a}+\underline{0}%
=\underline{0}+\mathbf{a}=\mathbf{a}$ for all $\mathbf{a}\in K\left[  \left[
x\right]  \right]  $.

\item \emph{Subtraction undoes addition:} Let $\mathbf{a},\mathbf{b}%
,\mathbf{c}\in K\left[  \left[  x\right]  \right]  $. We have $\mathbf{a}%
+\mathbf{b}=\mathbf{c}$ if and only if $\mathbf{a}=\mathbf{c}-\mathbf{b}$.

\item \emph{Commutativity of multiplication:} We have $\mathbf{ab}%
=\mathbf{ba}$ for all $\mathbf{a},\mathbf{b}\in K\left[  \left[  x\right]
\right]  $.

\item \emph{Associativity of multiplication:} We have $\mathbf{a}\left(
\mathbf{bc}\right)  =\left(  \mathbf{ab}\right)  \mathbf{c}$ for all
$\mathbf{a},\mathbf{b},\mathbf{c}\in K\left[  \left[  x\right]  \right]  $.

\item \emph{Distributivity:} We have%
\[
\mathbf{a}\left(  \mathbf{b}+\mathbf{c}\right)  =\mathbf{ab}+\mathbf{ac}%
\ \ \ \ \ \ \ \ \ \ \text{and}\ \ \ \ \ \ \ \ \ \ \left(  \mathbf{a}%
+\mathbf{b}\right)  \mathbf{c}=\mathbf{ac}+\mathbf{bc}%
\]
for all $\mathbf{a},\mathbf{b},\mathbf{c}\in K\left[  \left[  x\right]
\right]  $.

\item \emph{Neutrality of one:} We have $\mathbf{a}\cdot\underline{1}%
=\underline{1}\cdot\mathbf{a}=\mathbf{a}$ for all $\mathbf{a}\in K\left[
\left[  x\right]  \right]  $.

\item \emph{Annihilation:} We have $\mathbf{a}\cdot\underline{0}%
=\underline{0}\cdot\mathbf{a}=\underline{0}$ for all $\mathbf{a}\in K\left[
\left[  x\right]  \right]  $.
\end{enumerate}

\textbf{(b)} Furthermore, $K\left[  \left[  x\right]  \right]  $ is a
$K$-module (with its scaling being the map that sends each $\left(
\lambda,\mathbf{a}\right)  \in K\times K\left[  \left[  x\right]  \right]  $
to the FPS $\lambda\mathbf{a}$ defined in Definition \ref{def.fps.ops}
\textbf{(c)}). Its zero vector is $\underline{0}$. Concretely, this means that:

\begin{enumerate}
\item[10.] \emph{Associativity of scaling:} We have $\lambda\left(
\mu\mathbf{a}\right)  =\left(  \lambda\mu\right)  \mathbf{a}$ for all
$\lambda,\mu\in K$ and $\mathbf{a}\in K\left[  \left[  x\right]  \right]  $.

\item[11.] \emph{Left distributivity:} We have $\lambda\left(  \mathbf{a}%
+\mathbf{b}\right)  =\lambda\mathbf{a}+\lambda\mathbf{b}$ for all $\lambda\in
K$ and $\mathbf{a},\mathbf{b}\in K\left[  \left[  x\right]  \right]  $.

\item[12.] \emph{Right distributivity:} We have $\left(  \lambda+\mu\right)
\mathbf{a}=\lambda\mathbf{a}+\mu\mathbf{a}$ for all $\lambda,\mu\in K$ and
$\mathbf{a}\in K\left[  \left[  x\right]  \right]  $.

\item[13.] \emph{Neutrality of one:} We have $1\mathbf{a}=\mathbf{a}$ for all
$\mathbf{a}\in K\left[  \left[  x\right]  \right]  $.

\item[14.] \emph{Left annihilation:} We have $0\mathbf{a}=\underline{0}$ for
all $\mathbf{a}\in K\left[  \left[  x\right]  \right]  $.

\item[15.] \emph{Right annihilation:} We have $\lambda\underline{0}%
=\underline{0}$ for all $\lambda\in K$.
\end{enumerate}

(Also, some of the facts from part \textbf{(a)} are included in this
statement.) \medskip

\textbf{(c)} We have $\lambda\left(  \mathbf{a}\cdot\mathbf{b}\right)
=\left(  \lambda\mathbf{a}\right)  \cdot\mathbf{b}=\mathbf{a}\cdot\left(
\lambda\mathbf{b}\right)  $ for all $\lambda\in K$ and $\mathbf{a}%
,\mathbf{b}\in K\left[  \left[  x\right]  \right]  $. \medskip

\textbf{(d)} Finally, we have $\lambda\mathbf{a}=\underline{\lambda}%
\cdot\mathbf{a}$ for all $\lambda\in K$ and $\mathbf{a}\in K\left[  \left[
x\right]  \right]  $.
\end{theorem}

Theorem \ref{thm.fps.ring} allows us to calculate with FPSs as we do with
numbers, at least as far as the operations $+$, $-$ and $\cdot$ are concerned.
Hence, e.g., we know that:

\begin{itemize}
\item Sums and products in $K\left[  \left[  x\right]  \right]  $ need no
parentheses and do not depend on the order of addends/factors. For example,
for any $\mathbf{a},\mathbf{b},\mathbf{c},\mathbf{d}\in K\left[  \left[
x\right]  \right]  $, we have $\left(  \left(  \mathbf{ab}\right)
\mathbf{c}\right)  \mathbf{d}=\mathbf{a}\left(  \left(  \mathbf{bc}\right)
\mathbf{d}\right)  =\left(  \mathbf{ab}\right)  \left(  \mathbf{cd}\right)  $,
so that we can write $\mathbf{abcd}$ for each of these products; and moreover,
we have $\mathbf{abcd}=\mathbf{bdac}=\mathbf{dacb}$.

\item Finite sums and products (such as $\sum_{i=1}^{k}\mathbf{a}_{i}$ or
$\sum_{i\in I}\mathbf{a}_{i}$ or $\prod_{i=1}^{k}\mathbf{a}_{i}$ or
$\prod_{i\in I}\mathbf{a}_{i}$, where $I$ is a finite set) make sense and
behave as one would expect.

\item Powers exist: that is, you can take $\mathbf{a}^{n}$ for each FPS
$\mathbf{a}$ and each $n\in\mathbb{N}$.

\item Standard rules hold: e.g., we have $\mathbf{a}^{n+m}=\mathbf{a}%
^{n}\mathbf{a}^{m}$ and $\left(  \mathbf{ab}\right)  ^{n}=\mathbf{a}%
^{n}\mathbf{b}^{n}$ for any $\mathbf{a},\mathbf{b}\in K\left[  \left[
x\right]  \right]  $ and any $n,m\in\mathbb{N}$.

\item The binomial theorem holds: For any $\mathbf{a},\mathbf{b}\in K\left[
\left[  x\right]  \right]  $ and any $n\in\mathbb{N}$, we have%
\[
\left(  \mathbf{a}+\mathbf{b}\right)  ^{n}=\sum_{k=0}^{n}\dbinom{n}%
{k}\mathbf{a}^{k}\mathbf{b}^{n-k}.
\]

\end{itemize}

Before we prove Theorem \ref{thm.fps.ring}, we introduce one more notation:

\begin{definition}
\label{def.fps.coeff}If $n\in\mathbb{N}$, and if $\mathbf{a}=\left(
a_{0},a_{1},a_{2},\ldots\right)  \in K\left[  \left[  x\right]  \right]  $ is
an FPS, then we define an element $\left[  x^{n}\right]  \mathbf{a}\in K$ by%
\[
\left[  x^{n}\right]  \mathbf{a}:=a_{n}.
\]
This is called the \emph{coefficient of }$x^{n}$\emph{ in }$\mathbf{a}$, or
the $n$\emph{-th coefficient} of $\mathbf{a}$, or the $x^{n}$%
\emph{-coefficient} of $\mathbf{a}$.
\end{definition}

Thus, the definition of the sum of two FPSs (Definition \ref{def.fps.ops}
\textbf{(a)}) rewrites as follows: For any $\mathbf{a},\mathbf{b}\in K\left[
\left[  x\right]  \right]  $ and any $n\in\mathbb{N}$, we have
\begin{equation}
\left[  x^{n}\right]  \left(  \mathbf{a}+\mathbf{b}\right)  =\left[
x^{n}\right]  \mathbf{a}+\left[  x^{n}\right]  \mathbf{b}.
\label{pf.thm.fps.ring.xn(a+b)=}%
\end{equation}
Similarly, for any $\mathbf{a},\mathbf{b}\in K\left[  \left[  x\right]
\right]  $ and any $n\in\mathbb{N}$, we have
\begin{equation}
\left[  x^{n}\right]  \left(  \mathbf{a}-\mathbf{b}\right)  =\left[
x^{n}\right]  \mathbf{a}-\left[  x^{n}\right]  \mathbf{b}.
\label{pf.thm.fps.ring.xn(a-b)=}%
\end{equation}
Meanwhile, the definition of the product of two FPSs (Definition
\ref{def.fps.ops} \textbf{(d)}) rewrites as follows: For any $\mathbf{a}%
,\mathbf{b}\in K\left[  \left[  x\right]  \right]  $ and any $n\in\mathbb{N}$,
we have
\begin{align}
&  \left[  x^{n}\right]  \left(  \mathbf{ab}\right) \nonumber\\
&  =\left[  x^{0}\right]  \mathbf{a}\cdot\left[  x^{n}\right]  \mathbf{b}%
+\left[  x^{1}\right]  \mathbf{a}\cdot\left[  x^{n-1}\right]  \mathbf{b}%
+\left[  x^{2}\right]  \mathbf{a}\cdot\left[  x^{n-2}\right]  \mathbf{b}%
+\cdots+\left[  x^{n}\right]  \mathbf{a}\cdot\left[  x^{0}\right]
\mathbf{b}\nonumber\\
&  =\sum_{i=0}^{n}\left[  x^{i}\right]  \mathbf{a}\cdot\left[  x^{n-i}\right]
\mathbf{b}\label{pf.thm.fps.ring.xn(ab)=2}\\
&  =\sum_{j=0}^{n}\left[  x^{n-j}\right]  \mathbf{a}\cdot\left[  x^{j}\right]
\mathbf{b} \label{pf.thm.fps.ring.xn(ab)=3}%
\end{align}
(here, we have substituted $n-j$ for $i$ in the sum). For $n=0$, this equality
simplifies to%
\begin{equation}
\left[  x^{0}\right]  \left(  \mathbf{ab}\right)  =\left[  x^{0}\right]
\mathbf{a}\cdot\left[  x^{0}\right]  \mathbf{b}.
\label{pf.thm.fps.ring.x0(ab)=}%
\end{equation}
In other words, when we multiply two FPSs, their constant terms get
multiplied. Here and in the following, the \emph{constant term} of an FPS
$\mathbf{a}\in K\left[  \left[  x\right]  \right]  $ is defined to be its
$0$-th coefficient $\left[  x^{0}\right]  \mathbf{a}$.

Finally, Definition \ref{def.fps.ops} \textbf{(c)} rewrites as follows: For
any $\lambda\in K$ and $\mathbf{a}\in K\left[  \left[  x\right]  \right]  $
and any $n\in\mathbb{N}$, we have%
\begin{equation}
\left[  x^{n}\right]  \left(  \lambda\mathbf{a}\right)  =\lambda\cdot\left[
x^{n}\right]  \mathbf{a}. \label{pf.thm.fps.ring.xn(la)=}%
\end{equation}


\begin{proof}
[Proof of Theorem \ref{thm.fps.ring}.]Most parts of Theorem \ref{thm.fps.ring}
are straightforward to verify. Let us check the associativity of multiplication.

Let $\mathbf{a},\mathbf{b},\mathbf{c}\in K\left[  \left[  x\right]  \right]
$. We must prove that $\mathbf{a}\left(  \mathbf{bc}\right)  =\left(
\mathbf{ab}\right)  \mathbf{c}$. Let $n\in\mathbb{N}$. Consider the two
equalities%
\begin{align*}
\left[  x^{n}\right]  \left(  \left(  \mathbf{ab}\right)  \mathbf{c}\right)
&  =\sum_{j=0}^{n}\underbrace{\left[  x^{n-j}\right]  \left(  \mathbf{ab}%
\right)  }_{\substack{=\sum_{i=0}^{n-j}\left[  x^{i}\right]  \mathbf{a}%
\cdot\left[  x^{n-j-i}\right]  \mathbf{b}\\\text{(by
(\ref{pf.thm.fps.ring.xn(ab)=2}),}\\\text{applied to }n-j\text{ instead of
}n\text{)}}}\cdot\left[  x^{j}\right]  \mathbf{c}\\
&  \ \ \ \ \ \ \ \ \ \ \ \ \ \ \ \ \ \ \ \ \left(
\begin{array}
[c]{c}%
\text{by (\ref{pf.thm.fps.ring.xn(ab)=3}), applied to }\mathbf{ab}\text{ and
}\mathbf{c}\\
\text{instead of }\mathbf{a}\text{ and }\mathbf{b}%
\end{array}
\right) \\
&  =\sum_{j=0}^{n}\left(  \sum_{i=0}^{n-j}\left[  x^{i}\right]  \mathbf{a}%
\cdot\left[  x^{n-j-i}\right]  \mathbf{b}\right)  \cdot\left[  x^{j}\right]
\mathbf{c}\\
&  =\sum_{j=0}^{n}\ \ \sum_{i=0}^{n-j}\left[  x^{i}\right]  \mathbf{a}%
\cdot\left[  x^{n-j-i}\right]  \mathbf{b}\cdot\left[  x^{j}\right]  \mathbf{c}%
\end{align*}
and%
\begin{align*}
\left[  x^{n}\right]  \left(  \mathbf{a}\left(  \mathbf{bc}\right)  \right)
&  =\sum_{i=0}^{n}\left[  x^{i}\right]  \mathbf{a}\cdot\underbrace{\left[
x^{n-i}\right]  \left(  \mathbf{bc}\right)  }_{\substack{=\sum_{j=0}%
^{n-i}\left[  x^{n-i-j}\right]  \mathbf{b}\cdot\left[  x^{j}\right]
\mathbf{c}\\\text{(by (\ref{pf.thm.fps.ring.xn(ab)=3}), applied to
}n-i-j\text{, }\mathbf{b}\text{ and }\mathbf{c}\\\text{instead of }n\text{,
}\mathbf{a}\text{ and }\mathbf{b}\text{)}}}\\
&  \ \ \ \ \ \ \ \ \ \ \ \ \ \ \ \ \ \ \ \ \left(  \text{by
(\ref{pf.thm.fps.ring.xn(ab)=2}), applied to }\mathbf{bc}\text{ instead of
}\mathbf{b}\right) \\
&  =\sum_{i=0}^{n}\left[  x^{i}\right]  \mathbf{a}\cdot\left(  \sum
_{j=0}^{n-i}\left[  x^{n-i-j}\right]  \mathbf{b}\cdot\left[  x^{j}\right]
\mathbf{c}\right) \\
&  =\sum_{i=0}^{n}\ \ \sum_{j=0}^{n-i}\left[  x^{i}\right]  \mathbf{a}%
\cdot\left[  x^{n-i-j}\right]  \mathbf{b}\cdot\left[  x^{j}\right]
\mathbf{c}.
\end{align*}
The right hand sides of these two equalities are equal, since\footnote{The
first equality we are about to state is an equality of summation signs. Such
an equality means that whatever you put inside the summation signs, they
produce equal results. For example, $\sum_{i\in\left\{  1,2,3\right\}  }%
=\sum_{i=1}^{3}$ and $\sum_{i\in\mathbb{N}}=\sum_{i\geq0}$.}%
\[
\sum_{j=0}^{n}\ \ \sum_{i=0}^{n-j}=\sum_{\substack{\left(  i,j\right)
\in\mathbb{N}^{2};\\i+j\leq n}}=\sum_{i=0}^{n}\ \ \sum_{j=0}^{n-i}%
\]
and $n-j-i=n-i-j$. Thus, their left hand sides are equal as well. In other
words,%
\[
\left[  x^{n}\right]  \left(  \left(  \mathbf{ab}\right)  \mathbf{c}\right)
=\left[  x^{n}\right]  \left(  \mathbf{a}\left(  \mathbf{bc}\right)  \right)
.
\]


Now, forget that we fixed $n$. We thus have shown that $\left[  x^{n}\right]
\left(  \left(  \mathbf{ab}\right)  \mathbf{c}\right)  =\left[  x^{n}\right]
\left(  \mathbf{a}\left(  \mathbf{bc}\right)  \right)  $ for each
$n\in\mathbb{N}$. In other words, each entry of $\left(  \mathbf{ab}\right)
\mathbf{c}$ equals the corresponding entry of $\mathbf{a}\left(
\mathbf{bc}\right)  $. This entails $\left(  \mathbf{ab}\right)
\mathbf{c}=\mathbf{a}\left(  \mathbf{bc}\right)  $ (since a FPS is just the
sequence of its entries). In other words, $\mathbf{a}\left(  \mathbf{bc}%
\right)  =\left(  \mathbf{ab}\right)  \mathbf{c}$. This concludes the proof of
associativity of multiplication.

The remaining claims of Theorem \ref{thm.fps.ring} are
LTTR\footnote{\textquotedblleft LTTR\textquotedblright\ means
\textquotedblleft left to the reader\textquotedblright.} (their proofs follow
the same pattern, but are easier to execute).
\end{proof}

Since $K\left[  \left[  x\right]  \right]  $ is a commutative ring, any finite
sum of FPSs is well-defined. Sometimes, however, infinite sums of FPSs make
sense as well: for example, it stands to reason that%
\begin{align}
&  \ \ \ \ \left(  1,1,1,1,\ldots\right) \nonumber\\
&  +\left(  0,1,1,1,\ldots\right) \nonumber\\
&  +\left(  0,0,1,1,\ldots\right) \nonumber\\
&  +\left(  0,0,0,1,\ldots\right) \nonumber\\
&  +\cdots\nonumber\\
&  =\left(  1,2,3,4,\ldots\right)  , \label{eq.fps.summable.exa1}%
\end{align}
because FPSs are added entrywise. Let us rigorously define such sums. First,
we define \textquotedblleft essentially finite\textquotedblright\ sums of
elements of $K$:

\begin{definition}
\label{def.infsum.essfin}\textbf{(a)} A family $\left(  a_{i}\right)  _{i\in
I}\in K^{I}$ of elements of $K$ is said to be \emph{essentially finite} if all
but finitely many $i\in I$ satisfy $a_{i}=0$ (in other words, if the set
$\left\{  i\in I\ \mid\ a_{i}\neq0\right\}  $ is finite). \medskip

\textbf{(b)} Let $\left(  a_{i}\right)  _{i\in I}\in K^{I}$ be an essentially
finite family of elements of $K$. Then, the infinite sum $\sum_{i\in I}a_{i}$
is defined to equal the finite sum $\sum_{\substack{i\in I;\\a_{i}\neq0}%
}a_{i}$. Such an infinite sum is said to be \emph{essentially finite}.
\end{definition}

For example, the family $\left(  \left\lfloor \dfrac{5}{2^{n}}\right\rfloor
\right)  _{n\in\mathbb{N}}$ of integers is essentially finite\footnote{Here
and in the following, the notation $\left\lfloor x\right\rfloor $ (where $x$
is a real number) stands for the largest integer that is $\leq x$. For
instance, $\left\lfloor \pi\right\rfloor =3$ and $\left\lfloor 3\right\rfloor
=3$.}, and its sum is%
\begin{align*}
\sum_{n\in\mathbb{N}}\left\lfloor \dfrac{5}{2^{n}}\right\rfloor  &
=\left\lfloor \dfrac{5}{2^{0}}\right\rfloor +\left\lfloor \dfrac{5}{2^{1}%
}\right\rfloor +\left\lfloor \dfrac{5}{2^{2}}\right\rfloor +\left\lfloor
\dfrac{5}{2^{3}}\right\rfloor +\left\lfloor \dfrac{5}{2^{4}}\right\rfloor
+\cdots\\
&  =5+2+1+\underbrace{0+0+0+\cdots}_{\text{throw these away}}=5+2+1=8.
\end{align*}


Note the following:

\begin{itemize}
\item A family $\left(  a_{i}\right)  _{i\in I}\in K^{I}$ is always
essentially finite if $I$ is finite; thus, essentially finite families are a
wider class than finite families.

\item Any essentially finite sum of real or complex numbers is convergent in
the sense of analysis, but the converse is not true; for instance, the
infinite sum $\sum_{n\in\mathbb{N}}\dfrac{1}{2^{n}}=\dfrac{1}{1}+\dfrac{1}%
{2}+\dfrac{1}{4}+\dfrac{1}{8}+\cdots$ is convergent in the sense of analysis,
but not essentially finite. Essential finiteness is a crude algebraic
imitation of convergence. One of its advantage is that it works for sums of
elements of any commutative ring, not just of $\mathbb{R}$ or $\mathbb{C}$.
\end{itemize}

So the idea behind Definition \ref{def.infsum.essfin} \textbf{(b)} is that
addends that equal $0$ can be discarded in a sum, even when there are
infinitely many of them. \medskip

Sums of essentially finite families satisfy the usual rules for sums (such as
the breaking-apart rule $\sum_{i\in S}a_{s}=\sum_{i\in X}a_{s}+\sum_{i\in
Y}a_{s}$ when a set $S$ is the union of two disjoint sets $X$ and $Y$). See
\cite[\S 2.14.15]{detnotes} for details\footnote{Note that \cite[\S 2.14.15]%
{detnotes} uses the words \textquotedblleft\emph{finitely supported}%
\textquotedblright\ instead of \textquotedblleft essentially
finite\textquotedblright.}. There is only one \textbf{caveat}: Interchange of
summation signs (e.g., replacing $\sum_{i\in I}\ \ \sum_{j\in J}a_{i,j}$ by
$\sum_{j\in J}\ \ \sum_{i\in I}a_{i,j}$) works only if the family $\left(
a_{i,j}\right)  _{\left(  i,j\right)  \in I\times J}$ is essentially finite
(i.e., all but finitely many \textbf{pairs} $\left(  i,j\right)  \in I\times
J$ satisfy $a_{i,j}=0$); it does not suffice that the sums $\sum_{i\in
I}\ \ \sum_{j\in J}a_{i,j}$ and $\sum_{j\in J}\ \ \sum_{i\in I}a_{i,j}$
themselves are essentially finite (i.e., that the families $\left(
a_{i,j}\right)  _{j\in J}$ for all $i\in I$, the families $\left(
a_{i,j}\right)  _{i\in I}$ for all $j\in J$, and the families $\left(
\sum_{j\in J}a_{i,j}\right)  _{i\in I}$ and $\left(  \sum_{i\in I}%
a_{i,j}\right)  _{j\in J}$ are essentially finite).

\begin{fineprint}
For a counterexample, consider the family $\left(  a_{i,j}\right)  _{\left(
i,j\right)  \in I\times J}$ of integers with $I=\left\{  1,2,3,\ldots\right\}
$ and $J=\left\{  1,2,3,\ldots\right\}  $, where $a_{i,j}$ is given by the
following table:%
\[%
\begin{tabular}
[c]{|c||c|c|c|c|c|c|}\hline
$a_{i,j}$ & $j=1$ & $j=2$ & $j=3$ & $j=4$ & $j=5$ & $\cdots$\\\hline\hline
$i=1$ & $1$ & $-1$ &  &  &  & $\cdots$\\\hline
$i=2$ & $\phantom{-1}$ & $1$ & $-1$ &  &  & $\cdots$\\\hline
$i=3$ &  &  & $1$ & $-1$ &  & $\cdots$\\\hline
$i=4$ &  &  &  & $1$ & $-1$ & $\cdots$\\\hline
$i=5$ &  &  &  &  & $1$ & $\cdots$\\\hline
$\vdots$ & $\vdots$ & $\vdots$ & $\vdots$ & $\vdots$ & $\vdots$ & $\ddots
$\\\hline
\end{tabular}
\ \ \ \ \ \
\]
(where all the entries in the empty cells are $0$). For this family $\left(
a_{i,j}\right)  _{\left(  i,j\right)  \in I\times J}$, both sums $\sum_{i\in
I}\ \ \sum_{j\in J}a_{i,j}$ and $\sum_{j\in J}\ \ \sum_{i\in I}a_{i,j}$ are
essentially finite, but they are not equal (indeed, the former sum is
$\sum_{i\in I}\ \ \underbrace{\sum_{j\in J}a_{i,j}}_{=0}=\sum_{i\in I}0=0$,
whereas the latter sum is $\sum_{j\in J}\ \ \sum_{i\in I}a_{i,j}%
=\underbrace{\sum_{i\in I}a_{i,1}}_{=1}+\sum_{j>1}\ \ \underbrace{\sum_{i\in
I}a_{i,j}}_{=0}=1+\sum_{j>1}0=1$). And indeed, this family $\left(
a_{i,j}\right)  _{\left(  i,j\right)  \in I\times J}$ is not essentially finite.
\end{fineprint}

\bigskip

We have now made sense of infinite sums of elements of $K$ when all but
finitely many addends are $0$. Of course, we can do the same for $K\left[
\left[  x\right]  \right]  $ instead of $K$ (since $K\left[  \left[  x\right]
\right]  $, too, is a commutative ring). However, this does not help make
sense of the sum on the left hand side of (\ref{eq.fps.summable.exa1}),
because this sum is not essentially finite (it is a sum of infinitely many
nonzero FPSs). Thus, for sums of FPSs, we need a weaker version of essential
finiteness. Here is its definition:

\Needspace{19pc}

\begin{definition}
\label{def.fps.summable}A (possibly infinite) family $\left(  \mathbf{a}%
_{i}\right)  _{i\in I}$ of FPSs is said to be \emph{summable} (or
\emph{entrywise essentially finite}) if%
\[
\text{for each }n\in\mathbb{N}\text{, all but finitely many }i\in I\text{
satisfy }\left[  x^{n}\right]  \mathbf{a}_{i}=0.
\]
In this case, the sum $\sum_{i\in I}\mathbf{a}_{i}$ is defined to be the FPS
with%
\begin{equation}
\left[  x^{n}\right]  \left(  \sum_{i\in I}\mathbf{a}_{i}\right)
=\underbrace{\sum_{i\in I}\left[  x^{n}\right]  \mathbf{a}_{i}}%
_{\substack{\text{an essentially}\\\text{finite sum}}%
}\ \ \ \ \ \ \ \ \ \ \text{for all }n\in\mathbb{N}\text{.}
\label{eq.def.fps.summable.sum}%
\end{equation}

\end{definition}

\begin{remark}
The condition \textquotedblleft all but finitely many $i\in I$ satisfy
$\left[  x^{n}\right]  \mathbf{a}_{i}=0$\textquotedblright\ in Definition
\ref{def.fps.summable} is \textbf{not} equivalent to \textquotedblleft
infinitely many $i\in I$ satisfy $\left[  x^{n}\right]  \mathbf{a}_{i}%
=0$\textquotedblright.
\end{remark}

Any essentially finite family of FPSs is summable, but the converse is not
generally the case.

\begin{example}
Let us see how Definition \ref{def.fps.summable} justifies
(\ref{eq.fps.summable.exa1}). Consider the family $\left(  \mathbf{a}%
_{i}\right)  _{i\in\mathbb{N}}\in K\left[  \left[  x\right]  \right]
^{\mathbb{N}}$ of FPSs, where%
\[
\mathbf{a}_{i}:=\left(  \underbrace{0,0,\ldots,0}_{i\text{ zeroes}%
},1,1,1,\ldots\right)  \ \ \ \ \ \ \ \ \ \ \text{for each }i\in\mathbb{N}.
\]
The left hand side of (\ref{eq.fps.summable.exa1}) is precisely $\mathbf{a}%
_{0}+\mathbf{a}_{1}+\mathbf{a}_{2}+\cdots=\sum_{i\in\mathbb{N}}\mathbf{a}_{i}%
$, so let us check that the family $\left(  \mathbf{a}_{i}\right)
_{i\in\mathbb{N}}$ is summable and that its sum $\sum_{i\in\mathbb{N}%
}\mathbf{a}_{i}$ really equals the right hand side of
(\ref{eq.fps.summable.exa1}).

For each $n\in\mathbb{N}$, all but finitely many $i\in\mathbb{N}$ satisfy
$\left[  x^{n}\right]  \mathbf{a}_{i}=0$ (indeed, all $i>n$ satisfy this
equality). Thus, the family $\left(  \mathbf{a}_{i}\right)  _{i\in\mathbb{N}}$
is summable. For each $n\in\mathbb{N}$, we have%
\begin{align*}
\left[  x^{n}\right]  \left(  \sum_{i\in\mathbb{N}}\mathbf{a}_{i}\right)   &
=\sum_{i\in\mathbb{N}}\left[  x^{n}\right]  \mathbf{a}_{i}%
\ \ \ \ \ \ \ \ \ \ \left(  \text{by (\ref{eq.def.fps.summable.sum})}\right)
\\
&  =\sum_{i=0}^{n}\underbrace{\left[  x^{n}\right]  \mathbf{a}_{i}%
}_{\substack{=1\\\text{(since }i\leq n\text{)}}}+\sum_{i>n}\underbrace{\left[
x^{n}\right]  \mathbf{a}_{i}}_{\substack{=0\\\text{(since }i>n\text{)}}%
}=\sum_{i=0}^{n}1+\underbrace{\sum_{i>n}0}_{=0}\\
&  =\sum_{i=0}^{n}1=n+1.
\end{align*}
Thus, $\sum_{i\in\mathbb{N}}\mathbf{a}_{i}=\left(  1,2,3,4,\ldots\right)  $.
This is precisely the right hand side of (\ref{eq.fps.summable.exa1}). Thus,
(\ref{eq.fps.summable.exa1}) has been justified rigorously.
\end{example}

You can think of summable infinite sums of FPSs as a crude algebraic imitate
of uniformly convergent infinite sums of holomorphic functions in complex
analysis. (However, the former are a much simpler concept than the latter. In
particular, complex analysis is completely unnecessary in their study.)

The following fact is nearly obvious:\footnote{A \emph{subfamily} of a family
$\left(  f_{i}\right)  _{i\in I}$ means a family of the form $\left(
f_{i}\right)  _{i\in J}$, where $J$ is a subset of $I$.}

\begin{proposition}
\label{prop.fps.summable.sub}Let $\left(  \mathbf{a}_{i}\right)  _{i\in I}$ be
a summable family of FPSs. Then, any subfamily of $\left(  \mathbf{a}%
_{i}\right)  _{i\in I}$ is summable as well.
\end{proposition}

\begin{proof}
[Proof of Proposition \ref{prop.fps.summable.sub}.]Let $J$ be a subset of $I$.
We must prove that the subfamily $\left(  \mathbf{a}_{i}\right)  _{i\in J}$ is summable.

Let $n\in\mathbb{N}$. Then, all but finitely many $i\in I$ satisfy $\left[
x^{n}\right]  \mathbf{a}_{i}=0$ (since the family $\left(  \mathbf{a}%
_{i}\right)  _{i\in I}$ is summable). Hence, all but finitely many $i\in J$
satisfy $\left[  x^{n}\right]  \mathbf{a}_{i}=0$ (since $J$ is a subset of
$I$). Since we have proved this for each $n\in\mathbb{N}$, we thus conclude
that the family $\left(  \mathbf{a}_{i}\right)  _{i\in J}$ is summable. This
proves Proposition \ref{prop.fps.summable.sub}.
\end{proof}

Just as with essentially finite families, we can work with summable sums of
FPSs \textquotedblleft as if they were finite\textquotedblright\ most of the time:

\begin{proposition}
\label{prop.fps.summable-sums-rule}Sums of summable families of FPSs satisfy
the usual rules for sums (such as the breaking-apart rule $\sum_{i\in S}%
a_{s}=\sum_{i\in X}a_{s}+\sum_{i\in Y}a_{s}$ when a set $S$ is the union of
two disjoint sets $X$ and $Y$). See \cite[Proposition 7.2.11]{19s} for
details. Again, the only \textbf{caveat} is about interchange of summation
signs: The equality
\[
\sum_{i\in I}\ \ \sum_{j\in J}\mathbf{a}_{i,j}=\sum_{j\in J}\ \ \sum_{i\in
I}\mathbf{a}_{i,j}%
\]
holds when the family $\left(  \mathbf{a}_{i,j}\right)  _{\left(  i,j\right)
\in I\times J}$ is summable (i.e., when for each $n\in\mathbb{N}$, all but
finitely many \textbf{pairs} $\left(  i,j\right)  \in I\times J$ satisfy
$\left[  x^{n}\right]  \mathbf{a}_{i,j}=0$); it does not generally hold if we
merely assume that the sums $\sum_{i\in I}\ \ \sum_{j\in J}\mathbf{a}_{i,j}$
and $\sum_{j\in J}\ \ \sum_{i\in I}\mathbf{a}_{i,j}$ are summable.
\end{proposition}

\begin{proof}
[Proof of Proposition \ref{prop.fps.summable-sums-rule}.]The proof is tedious
(as there are many rules to check), but fairly straightforward (the idea is
always to focus on a single coefficient, and then to reduce the infinite sums
to finite sums). For example, consider the \textquotedblleft discrete Fubini
rule\textquotedblright, which says that%
\begin{equation}
\sum_{i\in I}\ \ \sum_{j\in J}\mathbf{a}_{i,j}=\sum_{\left(  i,j\right)  \in
I\times J}\mathbf{a}_{i,j}=\sum_{j\in J}\ \ \sum_{i\in I}\mathbf{a}_{i,j}
\label{pf.prop.fps.summable-sums-rule.fub}%
\end{equation}
whenever $\left(  \mathbf{a}_{i,j}\right)  _{\left(  i,j\right)  \in I\times
J}$ is a summable family of FPSs. In order to prove this rule, we fix a
summable family $\left(  \mathbf{a}_{i,j}\right)  _{\left(  i,j\right)  \in
I\times J}$ of FPSs. It is easy to see that the families $\left(
\mathbf{a}_{i,j}\right)  _{j\in J}$ for all $i\in I$ are summable as well, as
are the families $\left(  \mathbf{a}_{i,j}\right)  _{i\in I}$ for all $j\in
J$, and the families $\left(  \sum_{j\in J}\mathbf{a}_{i,j}\right)  _{i\in I}$
and $\left(  \sum_{i\in I}\mathbf{a}_{i,j}\right)  _{j\in J}$. Hence, all sums
in (\ref{pf.prop.fps.summable-sums-rule.fub}) are well-defined. Now, in order
to prove (\ref{pf.prop.fps.summable-sums-rule.fub}), it suffices to check that%
\[
\left[  x^{n}\right]  \left(  \sum_{i\in I}\ \ \sum_{j\in J}\mathbf{a}%
_{i,j}\right)  =\left[  x^{n}\right]  \left(  \sum_{\left(  i,j\right)  \in
I\times J}\mathbf{a}_{i,j}\right)  =\left[  x^{n}\right]  \left(  \sum_{j\in
J}\ \ \sum_{i\in I}\mathbf{a}_{i,j}\right)
\]
for each $n\in\mathbb{N}$. Fix $n\in\mathbb{N}$; then, we have $\left[
x^{n}\right]  \left(  \mathbf{a}_{i,j}\right)  =0$ for all but finitely many
$\left(  i,j\right)  \in I\times J$ (since the family $\left(  \mathbf{a}%
_{i,j}\right)  _{\left(  i,j\right)  \in I\times J}$ is summable). That is,
the set of all pairs $\left(  i,j\right)  \in I\times J$ satisfying $\left[
x^{n}\right]  \left(  \mathbf{a}_{i,j}\right)  \neq0$ is finite. Hence, the
set $I^{\prime}:=\left\{  i\ \mid\ \left(  i,j\right)  \in I\times J\text{
with }\left[  x^{n}\right]  \left(  \mathbf{a}_{i,j}\right)  \neq0\right\}  $
of the first entries of all these pairs is also finite, and so is the set
$J^{\prime}:=\left\{  i\ \mid\ \left(  i,j\right)  \in I\times J\text{ with
}\left[  x^{n}\right]  \left(  \mathbf{a}_{i,j}\right)  \neq0\right\}  $ of
the second entries of all these pairs. Now, the definitions of $I^{\prime}$
and $J^{\prime}$ ensure that any pair $\left(  i,j\right)  \in I\times J$
satisfies $\left[  x^{n}\right]  \mathbf{a}_{i,j}=0$ unless $i\in I^{\prime}$
and $j\in J^{\prime}$. Hence, we easily obtain the three equalities%
\[
\left[  x^{n}\right]  \left(  \sum_{i\in I}\ \ \sum_{j\in J}\mathbf{a}%
_{i,j}\right)  =\sum_{i\in I}\ \ \sum_{j\in J}\left[  x^{n}\right]
\mathbf{a}_{i,j}=\sum_{i\in I^{\prime}}\ \ \sum_{j\in J^{\prime}}\left[
x^{n}\right]  \mathbf{a}_{i,j}%
\]
and%
\[
\left[  x^{n}\right]  \left(  \sum_{\left(  i,j\right)  \in I\times
J}\mathbf{a}_{i,j}\right)  =\sum_{\left(  i,j\right)  \in I\times J}\left[
x^{n}\right]  \mathbf{a}_{i,j}=\sum_{\left(  i,j\right)  \in I^{\prime}\times
J^{\prime}}\left[  x^{n}\right]  \mathbf{a}_{i,j}%
\]
and%
\[
\left[  x^{n}\right]  \left(  \sum_{j\in J}\ \ \sum_{i\in I}\mathbf{a}%
_{i,j}\right)  =\sum_{j\in J}\ \ \sum_{i\in I}\left[  x^{n}\right]
\mathbf{a}_{i,j}=\sum_{j\in J^{\prime}}\ \ \sum_{i\in I^{\prime}}\left[
x^{n}\right]  \mathbf{a}_{i,j}.
\]
However, the right hand sides of these three equalities are equal (since the
sums appearing in them are finite sums, and thus satisfy the usual rules for
sums). Thus, the left hand sides are equal, exactly as we needed to show. See
\cite[proof of Proposition 7.2.11]{19s} for more details of this proof.
Proving the other properties of sums is easier.
\end{proof}

A few conventions about infinite sums will be used rather often:

\begin{convention}
\label{conv.fps.infsum}\textbf{(a)} For any given integer $m\in\mathbb{Z}$,
the summation sign $\sum_{k\geq m}$ is to be understood as $\sum_{k\in\left\{
m,m+1,m+2,\ldots\right\}  }$. We also write $\sum_{k=m}^{\infty}$ for this
summation sign. \medskip

\textbf{(b)} For any given integer $m\in\mathbb{Z}$, the summation sign
$\sum_{k>m}$ is to be understood as $\sum_{k\in\left\{  m+1,m+2,m+3,\ldots
\right\}  }$. \medskip

\textbf{(c)} Let $I$ be a set, and let $\mathcal{A}\left(  i\right)  $ be a
logical statement for each $i\in I$. (For example, $I$ can be $\mathbb{N}$,
and $\mathcal{A}\left(  i\right)  $ can be the statement \textquotedblleft$i$
is odd\textquotedblright.) Then, the summation sign $\sum_{\substack{i\in
I;\\\mathcal{A}\left(  i\right)  }}$ is to be understood as $\sum
_{i\in\left\{  j\in I\ \mid\ \mathcal{A}\left(  j\right)  \right\}  }$. (For
example, the summation sign $\sum_{\substack{i\in\mathbb{N};\\i\text{ is odd}%
}}$ means $\sum_{i\in\left\{  j\in\mathbb{N}\ \mid\ j\text{ is odd}\right\}
}$, that is, a sum over all odd elements of $\mathbb{N}$.)
\end{convention}

We can now define the $x$ that figured so prominently in our informal
exploration of formal power series back in Section \ref{sec.gf.exas}:

\begin{definition}
\label{def.fps.x}Let $x$ denote the FPS $\left(  0,1,0,0,0,\ldots\right)  $.
In other words, let $x$ denote the FPS with $\left[  x^{1}\right]  x=1$ and
$\left[  x^{i}\right]  x=0$ for all $i\neq1$.
\end{definition}

The following simple lemma follows almost immediately from the definition of
multiplication of FPSs:

\begin{lemma}
\label{lem.fps.xa}Let $\mathbf{a}=\left(  a_{0},a_{1},a_{2},\ldots\right)  $
be an FPS. Then, $x\cdot\mathbf{a}=\left(  0,a_{0},a_{1},a_{2},\ldots\right)
$.
\end{lemma}

In other words, multiplying an FPS $\mathbf{a}$ by $x$ is tantamount to
inserting a $0$ at the front of $\mathbf{a}$ (and shifting all the previously
existing entries of $\mathbf{a}$ to the right by one position).

\begin{proof}
[Proof of Lemma \ref{lem.fps.xa}.]If $n$ is a positive integer, then%
\begin{align*}
\left[  x^{n}\right]  \left(  x\cdot\mathbf{a}\right)   &  =\sum_{i=0}%
^{n}\left[  x^{i}\right]  x\cdot\underbrace{\left[  x^{n-i}\right]
\mathbf{a}}_{\substack{=a_{n-i}\\\text{(since }\mathbf{a}=\left(  a_{0}%
,a_{1},a_{2},\ldots\right)  \text{)}}}\\
&  \ \ \ \ \ \ \ \ \ \ \ \ \ \ \ \ \ \ \ \ \left(  \text{by
(\ref{pf.thm.fps.ring.xn(ab)=2}), applied to }x\text{ and }\mathbf{a}\text{
instead of }\mathbf{a}\text{ and }\mathbf{b}\right) \\
&  =\sum_{i=0}^{n}\left[  x^{i}\right]  x\cdot a_{n-i}=\underbrace{\left[
x^{0}\right]  x}_{=0}\cdot\,a_{n-0}+\underbrace{\left[  x^{1}\right]  x}%
_{=1}\cdot\,a_{n-1}+\sum_{i=2}^{n}\underbrace{\left[  x^{i}\right]
x}_{\substack{=0\\\text{(since }i\geq2>1\text{)}}}\cdot\,a_{n-i}\\
&  \ \ \ \ \ \ \ \ \ \ \ \ \ \ \ \ \ \ \ \ \left(
\begin{array}
[c]{c}%
\text{here, we have split off the addends}\\
\text{for }i=0\text{ and }i=1\text{ from the sum}%
\end{array}
\right) \\
&  =\underbrace{0\cdot a_{n-0}}_{=0}+\,1\cdot a_{n-1}+\underbrace{\sum
_{i=2}^{n}0\cdot a_{n-i}}_{=0}=1\cdot a_{n-1}=a_{n-1}.
\end{align*}
A similar argument can be used for $n=0$ (except that now, the sum $\sum
_{i=0}^{n}\left[  x^{i}\right]  x\cdot a_{n-i}$ has no $\left[  x^{1}\right]
x\cdot a_{n-1}$ addend), and results in the conclusion that $\left[
x^{n}\right]  \left(  x\cdot\mathbf{a}\right)  =0$ in this case. Thus, for
each $n\in\mathbb{N}$, we have%
\[
\left[  x^{n}\right]  \left(  x\cdot\mathbf{a}\right)  =%
\begin{cases}
a_{n-1}, & \text{if }n>0;\\
0, & \text{if }n=0.
\end{cases}
\]
In other words, $x\cdot\mathbf{a}=\left(  0,a_{0},a_{1},a_{2},\ldots\right)
$. This proves Lemma \ref{lem.fps.xa}.
\end{proof}

Recall that $x^{k}=\underbrace{xx\cdots x}_{k\text{ times}}$ for each
$k\in\mathbb{N}$ (by the definition of powers in any commutative ring). In
particular, $x^{0}=\underline{1}$ (since $\underline{1}$ is the unity of the
ring $K\left[  \left[  x\right]  \right]  $). The following proposition
describes $x^{k}$ explicitly for each $k\in\mathbb{N}$:

\begin{proposition}
\label{prop.fps.xk}We have%
\[
x^{k}=\left(  \underbrace{0,0,\ldots,0}_{k\text{ times}},1,0,0,0,\ldots
\right)  \ \ \ \ \ \ \ \ \ \ \text{for each }k\in\mathbb{N}.
\]

\end{proposition}

\begin{proof}
[Proof of Proposition \ref{prop.fps.xk}.]We induct on $k$.

\textit{Induction base:} We have $x^{0}=\underline{1}=\left(  1,0,0,0,0,\ldots
\right)  =\left(  \underbrace{0,0,\ldots,0}_{0\text{ times}},1,0,0,0,\ldots
\right)  $. In other words, Proposition \ref{prop.fps.xk} holds for $k=0$.

\textit{Induction step:} Let $m\in\mathbb{N}$. Assume that Proposition
\ref{prop.fps.xk} holds for $k=m$. We must prove that Proposition
\ref{prop.fps.xk} holds for $k=m+1$.

We have $x^{m}=\left(  \underbrace{0,0,\ldots,0}_{m\text{ times}%
},1,0,0,0,\ldots\right)  $ (since Proposition \ref{prop.fps.xk} holds for
$k=m$). Thus, Lemma \ref{lem.fps.xa} (applied to $\mathbf{a}=x^{m}$ and
\newline$\left(  a_{0},a_{1},a_{2},\ldots\right)  =\left(
\underbrace{0,0,\ldots,0}_{m\text{ times}},1,0,0,0,\ldots\right)  $) yields%
\[
x\cdot x^{m}=\left(  0,\underbrace{0,0,\ldots,0}_{m\text{ times}%
},1,0,0,0,\ldots\right)  =\left(  \underbrace{0,0,\ldots,0}_{m+1\text{ times}%
},1,0,0,0,\ldots\right)  .
\]
In other words, $x^{m+1}=\left(  \underbrace{0,0,\ldots,0}_{m+1\text{ times}%
},1,0,0,0,\ldots\right)  $ (since $x\cdot x^{m}=x^{m+1}$). In other words,
Proposition \ref{prop.fps.xk} holds for $k=m+1$. This completes the induction
step, thus proving Proposition \ref{prop.fps.xk}.
\end{proof}

Finally, the following corollary allows us to rewrite any FPS $\left(
a_{0},a_{1},a_{2},\ldots\right)  $ in the familiar form $a_{0}+a_{1}%
x+a_{2}x^{2}+a_{3}x^{3}+\cdots$:

\begin{corollary}
\label{cor.fps.sumakxk}Any FPS $\left(  a_{0},a_{1},a_{2},\ldots\right)  \in
K\left[  \left[  x\right]  \right]  $ satisfies%
\[
\left(  a_{0},a_{1},a_{2},\ldots\right)  =a_{0}+a_{1}x+a_{2}x^{2}+a_{3}%
x^{3}+\cdots=\sum_{n\in\mathbb{N}}a_{n}x^{n}.
\]
In particular, the right hand side here is well-defined, i.e., the family
$\left(  a_{n}x^{n}\right)  _{n\in\mathbb{N}}$ is summable.
\end{corollary}

\begin{proof}
[Proof of Corollary \ref{cor.fps.sumakxk} (sketched).](See \cite[Corollary
7.2.16]{19s} for details.) By Proposition \ref{prop.fps.xk}, we have%
\begin{align*}
&  a_{0}+a_{1}x+a_{2}x^{2}+a_{3}x^{3}+\cdots\\
&  =\ \ \ \ a_{0}\left(  1,0,0,0,\ldots\right) \\
&  \ \ \ \ +a_{1}\left(  0,1,0,0,\ldots\right) \\
&  \ \ \ \ +a_{2}\left(  0,0,1,0,\ldots\right) \\
&  \ \ \ \ +a_{3}\left(  0,0,0,1,\ldots\right) \\
&  \ \ \ \ +\cdots\\
&  =\ \ \ \ \left(  a_{0},0,0,0,\ldots\right) \\
&  \ \ \ \ +\left(  0,a_{1},0,0,\ldots\right) \\
&  \ \ \ \ +\left(  0,0,a_{2},0,\ldots\right) \\
&  \ \ \ \ +\left(  0,0,0,a_{3},\ldots\right) \\
&  \ \ \ \ +\cdots\\
&  =\left(  a_{0},a_{1},a_{2},a_{3},\ldots\right)  .
\end{align*}
This proves Corollary \ref{cor.fps.sumakxk} (since $a_{0}+a_{1}x+a_{2}%
x^{2}+a_{3}x^{3}+\cdots=\sum_{n\in\mathbb{N}}a_{n}x^{n}$ holds for obvious reasons).
\end{proof}

So we have \textquotedblleft found\textquotedblright\ our $x$ and given a
rigorous justification for writing $\left(  a_{0},a_{1},a_{2},\ldots\right)  $
as $a_{0}+a_{1}x+a_{2}x^{2}+a_{3}x^{3}+\cdots$. Note that we did not use any
analysis (real or complex) in the process; in particular, we did not have to
worry about convergence (we did have to worry about summability, but this is
much simpler than convergence). It is easy to come up with FPSs that don't
converge when any nonzero real number is substituted for $x$ (for example,
$\sum_{n\in\mathbb{N}}n!x^{n}=1+x+2x^{2}+6x^{3}+24x^{4}+\cdots$ is such an
FPS); they are nevertheless completely legitimate FPSs.

We can now also answer the question \textquotedblleft what is a generating
function\textquotedblright, albeit the answer is somewhat anticlimactic:

\begin{definition}
Let $\left(  a_{0},a_{1},a_{2},\ldots\right)  $ be a sequence of elements of
$K$. Then, the \emph{(ordinary) generating function} of $\left(  a_{0}%
,a_{1},a_{2},\ldots\right)  $ will mean the FPS $\left(  a_{0},a_{1}%
,a_{2},\ldots\right)  =a_{0}+a_{1}x+a_{2}x^{2}+a_{3}x^{3}+\cdots$.
\end{definition}

\subsubsection{\label{subsec.gf.defs.cvi}The Chu--Vandermonde identity}

What we have done so far is sufficient to justify Example 3 from Section
\ref{sec.gf.exas}. Thus, let us record the result of Example 3 as a proposition:

\begin{proposition}
\label{prop.binom.vandermonde.NN}Let $a,b\in\mathbb{N}$, and let
$n\in\mathbb{N}$. Then,%
\begin{equation}
\dbinom{a+b}{n}=\sum_{k=0}^{n}\dbinom{a}{k}\dbinom{b}{n-k}.
\label{eq.prop.binom.vandermonde.NN.eq}%
\end{equation}

\end{proposition}

We have yet to justify Examples 1, 2 and 4; we shall do so later. First,
however, let us generalize Proposition \ref{prop.binom.vandermonde.NN} to
arbitrary numbers $a,b$ (as opposed to merely $a,b\in\mathbb{N}$). That is, we
shall prove the following:

\begin{theorem}
[\emph{Vandermonde convolution identity}, aka \emph{Chu--Vandermonde
identity}]\label{thm.binom.vandermonde.CC}Let $a,b\in\mathbb{C}$, and let
$n\in\mathbb{N}$. Then,%
\begin{equation}
\dbinom{a+b}{n}=\sum_{k=0}^{n}\dbinom{a}{k}\dbinom{b}{n-k}.
\label{eq.prop.binom.vandermonde.CC.eq}%
\end{equation}

\end{theorem}

(Actually, the \textquotedblleft Let $a,b\in\mathbb{C}$\textquotedblright\ in
Theorem \ref{thm.binom.vandermonde.CC} can be replaced by \textquotedblleft
Let $a,b$ be any numbers\textquotedblright, where \textquotedblleft
numbers\textquotedblright\ is understood appropriately\footnote{Namely, a
\textquotedblleft number\textquotedblright\ should here be viewed as an
element of a commutative $\mathbb{Q}$-algebra. This includes complex numbers,
polynomials over complex numbers, power series over complex numbers and even
commuting matrices with complex entries.}. We just wrote \textquotedblleft Let
$a,b\in\mathbb{C}$\textquotedblright\ for simplicity.)

To recall, back in Example 3, we proved Proposition
\ref{prop.binom.vandermonde.NN} by multiplying out the identity $\left(
1+x\right)  ^{a+b}=\left(  1+x\right)  ^{a}\left(  1+x\right)  ^{b}$ using the
binomial formula. If we wanted to extend this argument to arbitrary
$a,b\in\mathbb{C}$, then we would need to make sense of powers like $\left(
1+x\right)  ^{-3}$ or $\left(  1+x\right)  ^{1/2}$ or $\left(  1+x\right)
^{\pi}$. This is indeed possible (in fact, we will briefly outline this later
on), but there is a much shorter way.

In fact, there is a slick trick to automatically extend a claim like
(\ref{eq.prop.binom.vandermonde.NN.eq}) from nonnegative integers to complex
numbers. It is sometimes known as the \emph{polynomial identity trick}, and is
used a lot in algebra. The proof of Theorem \ref{thm.binom.vandermonde.CC}
that we shall sketch below should illustrate this trick; you can find more
details and further examples in \cite[\S 7.5.3]{20f}.

\begin{proof}
[Proof of Theorem \ref{thm.binom.vandermonde.CC} (sketched).]Fix
$n\in\mathbb{N}$ and $b\in\mathbb{N}$, but forget that $a$ was fixed. Then,
Proposition \ref{prop.binom.vandermonde.NN} (which we have already proved)
says that the equality%
\begin{equation}
\dbinom{a+b}{n}=\sum_{k=0}^{n}\dbinom{a}{k}\dbinom{b}{n-k}
\label{pf.thm.binom.vandermonde.CC.1}%
\end{equation}
holds for each $a\in\mathbb{N}$. However, both sides of this equality are
polynomials (more precisely, polynomial functions) in $a$: indeed,%
\begin{align*}
\dbinom{a+b}{n}  &  =\dfrac{\left(  a+b\right)  \left(  a+b-1\right)  \left(
a+b-2\right)  \cdots\left(  a+b-n+1\right)  }{n!}%
\ \ \ \ \ \ \ \ \ \ \text{and}\\
\sum_{k=0}^{n}\dbinom{a}{k}\dbinom{b}{n-k}  &  =\sum_{k=0}^{n}\dfrac{a\left(
a-1\right)  \cdots\left(  a-k+1\right)  }{k!}\dbinom{b}{n-k}.
\end{align*}
If two univariate polynomials $p$ and $q$ (with rational, real or complex
coefficients) are equal for each $a\in\mathbb{N}$ (that is: if we have
$p\left(  a\right)  =q\left(  a\right)  $ for each $a\in\mathbb{N}$), then
they must be identical (because two univariate polynomials that are equal at
infinitely many points must necessarily be identical\footnote{Quick reminder
on why this is true: If $p$ and $q$ are two univariate polynomials (with
rational, real or complex coefficients) that are equal at infinitely many
points (i.e., if there exist infinitely many numbers $z$ satisfying $p\left(
z\right)  =q\left(  z\right)  $), then $p=q$ (because the assumption entails
that the difference $p-q$ has infinitely many roots, but this entails $p-q=0$
and thus $p=q$). See \cite[Corollary 7.5.7]{20f} for this argument in more
detail.}). Hence, the two sides of (\ref{pf.thm.binom.vandermonde.CC.1}) must
be identical as polynomials in $a$. Thus, the equality
(\ref{pf.thm.binom.vandermonde.CC.1}) holds not only for each $a\in\mathbb{N}%
$, but also for each $a\in\mathbb{C}$.

Now, forget that $b$ was fixed. Instead, let us fix $a\in\mathbb{C}$. As we
just have proved, the equality (\ref{pf.thm.binom.vandermonde.CC.1}) holds for
each $b\in\mathbb{N}$. We want to show that it holds for each $b\in\mathbb{C}%
$. But this can be achieved by the same argument that we just used to extend
it from $a\in\mathbb{N}$ to $a\in\mathbb{C}$: We view both sides of the
equality as polynomials (but this time in $b$, not in $a$), and argue that
these polynomials must be identical because they are equal at infinitely many
points. The upshot is that the equality (\ref{pf.thm.binom.vandermonde.CC.1})
holds for all $a,b\in\mathbb{C}$; thus, Theorem \ref{thm.binom.vandermonde.CC}
is proven. (See \cite[proofs of Lemma 7.5.8 and Theorem 7.5.3]{20f} or
\cite[\S 2.17.3]{19s} for this proof in more detail. Alternatively, see
\cite[\S 3.3.2 and \S 3.3.3]{detnotes} for two other proofs.)
\end{proof}

\subsubsection{What next?}

Let us now return to our quest of justifying Examples 1, 2 and 4 from Section
\ref{sec.gf.exas}. In order to do so, we need to know

\begin{itemize}
\item what we can substitute into a FPS;

\item when and why can we divide FPSs by FPSs;

\item when and why can we take the square root of an FPS and solve a quadratic
equation using the quadratic formula.
\end{itemize}

So we need to do more. The following sections are devoted to this.
